\documentclass[aspectratio=169,10pt,english]{beamer}

\input{../../../_preambulo_beamer.tex}
\input{../../../_comandos_beamer.tex}

\selectlanguage{english}

\title{MA1001B - Statistical Modeling for Decision Making}
\subtitle{Lesson 36: Hypothesis Test for a Business Proportion}
\author{}
\institute{Tecnol\'ogico de Monterrey}
\date{}

\begin{document}

\begin{frame}[plain]
  \vspace{1.2cm}
  {\Large\bfseries MA1001B - Statistical Modeling for Decision Making\par}
  \vspace{0.7cm}
  {\large Lesson 36: Hypothesis Test for a Business Proportion\par}
  \vspace{0.7cm}
  {\color{TecRojo}\rule{\textwidth}{1.2pt}}
  \vspace{0.8cm}

  MA1001B Faculty\\[0.3cm]
  Term\\[0.3cm]
  Tecnol\'ogico de Monterrey
\end{frame}

\begin{frame}{The One-Proportion Question}
  \begin{alertblock}{Guiding question}
    How do we test a claimed late-flag rate $p=0.10$ from $x=28$ flags in
    $n=200$ lots?
  \end{alertblock}

  \vspace{0.6em}
  By the end of this lesson, you can:
  \begin{itemize}
    \item write $H_0:p=0.10$ versus $H_1:p>0.10$;
    \item check $np_0>5$ and $n(1-p_0)>5$;
    \item compute $z$ with $p_0$ in the standard error;
    \item explain why replacing $p_0$ with $\hat{p}$ is the wrong statistic.
  \end{itemize}

  \vfill
  \scriptsize All lots and late flags used today are synthetic. No real company data are used.
\end{frame}

\begin{frame}{A Claimed Rate and a Sample Count}
  \small
  A synthetic inbound-inspection process claims that 10 percent of lots are
  flagged late. The observed register is $x=28$ late flags in $n=200$ lots.

  \[
    \hat{p}=\frac{28}{200}=0.14,\qquad
    np_0=200\times0.10=20,\qquad
    n(1-p_0)=180.
  \]

  \begin{exampleblock}{Right-tailed hypotheses}
    $H_0:p=0.10$ versus $H_1:p>0.10$, with $\alpha=0.05$. The normal $z$
    approximation is used only after $np_0>5$ and $n(1-p_0)>5$.
  \end{exampleblock}
\end{frame}

\begin{frame}[fragile]{Step 1: Sample Proportion and Conditions}
  \begin{columns}[T]
    \begin{column}{0.42\textwidth}
      \small
      \begin{lstlisting}
phat = 28 / 200
np0 = n * p0
nq0 = n * (1 - p0)
      \end{lstlisting}

      \begin{block}{Verified output}
        $\hat{p}=0.140000$\\
        $np_0=20$, $n(1-p_0)=180$\\
        Conditions met: yes
      \end{block}
    \end{column}
    \begin{column}{0.58\textwidth}
      \centering
      \includegraphics[width=\textwidth]{../figures/proportion_test_01_sample_and_conditions.png}
      \vspace{0.2em}
      \scriptsize Hypothesized $p_0$ versus $\hat{p}$ from Step 1
    \end{column}
  \end{columns}
\end{frame}

\begin{frame}{The Test Standard Error Uses $p_0$}
  \small
  Under $H_0:p=p_0$ the sampling variance is $p_0(1-p_0)/n$, not
  $\hat{p}(1-\hat{p})/n$:
  \[
    z=\frac{\hat{p}-p_0}{\sqrt{p_0(1-p_0)/n}}
    =\frac{0.14-0.10}{0.021213}=1.885618.
  \]
  The right-tailed p-value is
  \[
    p=P(Z\ge 1.885618)=0.029673.
  \]
  Because $p<\alpha=0.05$ and $z>z^*=1.644854$, the test rejects $H_0$.
\end{frame}

\begin{frame}[fragile]{Step 2: $z$ and the Right-Tailed p-Value}
  \begin{columns}[T]
    \begin{column}{0.40\textwidth}
      \small
      \begin{lstlisting}
se = sqrt(p0*(1-p0)/n)
z = (phat - p0) / se
p = norm.sf(z)
      \end{lstlisting}

      \begin{block}{Verified output}
        $SE=0.021213$, $z=1.885618$\\
        $p=0.029673$, $z^*=1.644854$\\
        Decision: reject $H_0$
      \end{block}
    \end{column}
    \begin{column}{0.60\textwidth}
      \centering
      \includegraphics[width=\textwidth]{../figures/proportion_test_02_z_pvalue.png}
      \vspace{0.2em}
      \scriptsize Right-tailed $z$ test from Step 2
    \end{column}
  \end{columns}
\end{frame}

\begin{frame}{Step 3: $\hat{p}$ in the SE Is the Wrong Statistic}
  \begin{columns}[T]
    \begin{column}{0.42\textwidth}
      \small
      Correct SE uses $p_0$: $0.021213$, $z=1.885618$, $p=0.029673$, reject.

      \vspace{0.5em}
      Wrong SE uses $\hat{p}$: $0.024536$, $z=1.630278$, $p=0.051521$,
      do not reject.

      \vspace{0.5em}
      \textbf{Limit:} a test of $H_0:p=p_0$ must put $p_0$ in the standard
      error. Using $\hat{p}$ is a confidence-interval SE and can flip the
      $\alpha=0.05$ decision.
    \end{column}
    \begin{column}{0.58\textwidth}
      \centering
      \includegraphics[width=\textwidth]{../figures/proportion_test_03_wrong_se_limit.png}
      \vspace{0.2em}
      \scriptsize Correct versus wrong SE from Step 3
    \end{column}
  \end{columns}
\end{frame}

\begin{frame}{Concept Check}
  \begin{enumerate}
    \item Compute $\hat{p}$ from $x=28$ and $n=200$.
    \item Why is $np_0=20$ checked with $p_0=0.10$ rather than with $\hat{p}$?
    \item Write the $z$ formula that belongs in this test.
    \item Why can using $\hat{p}$ in the SE flip the decision at $\alpha=0.05$?
  \end{enumerate}

  \vspace{0.6em}
  \begin{block}{Connection to Lesson 37}
    The session can continue directly into a two-sample $t$ test for two
    independent means, where each group has its own $s$ and $n$.
  \end{block}

  \vfill
  \scriptsize Source: OpenStax, \textit{Introductory Business Statistics 2e},
  Chapter 9, Sections 9.3--9.4. CC BY-NC-SA.
\end{frame}

\appendix



\end{document}
